\section{Pointwise Truth Is Not Connection}
\label{sec:pointwise-truth-connection}

The separation principle has one classical part and one Cred-specific
part.  The classical part is the distinction between truth at one
evaluation and consequence across all evaluations.  If two interpreted
formulas both have value $1$ at the present valuation, this does not
make one a consequence of the other.  Equivalently, two functions may
satisfy $f(x_0)=g(x_0)=1$ at one point without satisfying $f=g$ or
$f\le g$ on the whole domain.  Classical model theory already uses
this distinction: consequence is a metalinguistic relation over all
models, not an object-language truth function.  This is not the novelty
of Cred.

The Cred-specific part is the treatment of conditional dependence.  A
truth value is not a connection.  In probability, the marginal values
$P(A)$ and $P(B)$ do not determine the joint value $P(A\cap B)$.  The
joint is additional dependence data.  Independence is the special case
\[
  P(A\cap B)=P(A)P(B),
\]
which yields $P(A\mid B)=P(A)$ whenever $P(B)>0$.  In that case the
evidence has no updating force for the conclusion.  Dependence appears
only when the supplied joint differs from the product joint.

Cred imports exactly this lesson at the scalar level.  The values of
$A$ and $B$ do not determine the conditional value.  One must supply a
joint value $j=\cred(A\land B)$, or equivalently a conditional candidate
$c=\cred(A\mid B)$ together with the evidence value $e=\cred(B)$, and
then require the chain-rule constraint
\[
  \Adm(c;j,e)\quad\Longleftrightarrow\quad c e=j.
\]
At positive evidence the relation is invertible in the usual direction
and gives $c=j/e$.  At zero evidence the inverse loses rank:
\[
  c\cdot 0=0
\]
for every $c$.  Thus impossible evidence is not a false premise that
licenses arbitrary conclusions.  It is a rank-deficient constraint.
The admissible set is the whole fiber $\Cond(0,0)=[0,1]$.

This also clarifies the relation to relevance logic.  Relevance logics
add proof-theoretic or semantic constraints on premise use, topic
sharing, or aboutness.  Cred does not need that extra layer to define
entailment: entailment remains an external consequence relation.
Instead, Cred needs supplied dependence data when the question is not
validity but connection.  The slogan is therefore:
\[
\boxed{\text{a formula has a value; a connection requires extra structure.}}
\]
The unique advantage claimed here is not that truth differs from
consequence.  Classical logic already knows that.  The advantage is a
small algebraic interface in which truth values, supplied dependence,
and conditional inference are kept separate, and in which zero evidence
propagates a fiber instead of manufacturing a vacuous truth value.
